\section{从竞争到协作：LLM Agent-Agent Cooperation 的机会与硬约束}

\subsection{为什么需要多 Agent 协作}
Brown 先从 reasoning model 的成功讲起：test-time compute 增加，性能持续上升。但串行 Chain-of-Thought 存在硬延迟边界，长时任务不可能无限线性等待。

\begin{knowledgebox}{协作侧动机}
\begin{itemize}
    \item \textbf{Latency}：并行采样/并行推理可换取墙钟时间。
    \item \textbf{Diversity}：不同模型/不同策略在子问题上有互补优势。
    \item \textbf{Routing}：任务可分发到更擅长的专家模型，类似工具调用。
\end{itemize}
\end{knowledgebox}

\subsection{Consensus 与 Best-of-N}
从串行推理转向并行协作后，最先需要建立的是一个朴素基线：同一问题多采样几次，究竟该按票数聚合，还是交给验证器挑选？这两种方法都没有真正解决开放式协商，却能把“多 Agent 是否有增益”拆成可测量的样本多样性、答案相关性与判别质量问题。
课程中比较了两类简单但高频的并行策略：
\begin{itemize}
    \item \textbf{Consensus}：多次采样后取多数答案，适合短答案任务；
    \item \textbf{Best-of-N}：多样本后由验证器选最优，适合可验证任务。
\end{itemize}

\begin{importantbox}{性能与效率权衡}
这类方法能降低延迟并提升上界，但通常 \textbf{计算效率更差}，且依赖答案可比较性或可验证性，不是通用解。
\end{importantbox}

\subsection{Diversity 与模型路由}
Brown 用 ``大数乘法应交给计算器'' 的比喻说明：最强模型不应吞掉所有任务。多 Agent 的一个高性价比形式，是把路由（routing）当作策略层，让查询进入最匹配专家。

\begin{warningbox}{把多 Agent 当银弹的风险}
过度堆叠 scaffold 往往得到 ``更复杂但更脆'' 的系统。若没有稳定通信协议、冲突仲裁与状态一致性，多 Agent 只会放大错误传播。
\end{warningbox}

\subsection{本章小结}
协作式多 Agent 的收益真实存在，但当前更多体现在可验证子任务和路由层。开放式长链协商仍是短板。

